%------------------------------------------------------------------------------%
\subsection{Problema de Basileia}
\label{subsec:problema-Basileia}

O \emph{Problema de Basileia}\index{Problema!Basileia} consiste em
encontrar a soma da Série de Inversos do Quadrado,
apresentada na Definição~\ref{def:serie-inverso-quadrado},
\[
\sum_{n=1}^\infty \frac{1}{\,n^2\,} = \frac{\,\pi^2}{6}    \;.
\]
A solução desse problema obtida por Euler em 1735 o tornou conhecido na comunidade
Matemática da época e é de grande importância
por estar relacionada ao estuda da Função Zeta de Riemann, $\zeta(z)$,
pois
\[
\zeta(2) = \sum_{n=1}^\infty \frac{1}{\,n^2\,}    \;.
\]
Entre as aplicações desse resultado está o cálculo da probabilidade de dois
números escolhidos ao acaso serem primos entre si.
A verificação dos cálculos realizados por Euler usa a decomposição do seno
em um produto infinito que depende no Teorema de Fatoração de Weierstrass
para funções holomorfas
em todo o plano complexo, evidentemente além do escopo dessa disciplina.
Dessa forma, vamos apresentar apenas os passos principais desse desenvolvimento
e em seguida apresentamos uma demonstração
publicada por T. M. Apostol~\cite{Apostol-Proof-Euler-Missed} em 1983
que envolve operações mais simples.

O desenvolvimento dos cálculos de Euler começa com a Série de Taylor da função seno
\[
\sin(x) = x
- \frac{x^3}{3!}
+ \frac{x^5}{5!}
- \frac{x^7}{7!}
+ \cdots
\]
dividindo ambos os lados por $x$ temos
\[
\frac{\sin(x)}{x} = 1
- \frac{x^2}{3!}
+ \frac{x^4}{5!}
- \frac{x^6}{7!}
+ \cdots
\]
Usando o Teorema de Fatoração de Weierstrass nessa função, podemos escrever
\[
\frac{\sin(x)}{x} = \prod_{n=1}^{\infty} \left(1-\frac{x^2}{\pi^2n^2}\right)
= \left(1-\frac{x^2}{ \pi^2}\right)
\left(1-\frac{x^2}{4\pi^2}\right)
\left(1-\frac{x^2}{9\pi^2}\right)
\cdots
\]
Igualando a Série de Taylor da função com sua fatoração temos
\begin{equation}
  \label{eq:fatoracao_seno}
  1
  - \frac{x^2}{3!}
  + \frac{x^4}{5!}
  - \frac{x^6}{7!}
  + \cdots
  = \left(1-\frac{x^2}{ \pi^2}\right)
  \left(1-\frac{x^2}{4\pi^2}\right)
  \left(1-\frac{x^2}{9\pi^2}\right)
  \cdots
\end{equation}
Como essa igualdade é válida sabemos que o coeficiente de cada potência
de $x$ é igual nas duas expressões.
Extraímos diretamente o coeficiente de $x^2$ da Série de Taylor
obtendo
\[
-\frac{1}{3!} = -\frac{1}{6}    \;.
\]
Multiplicando os termos do produto e coletando apenas os
coeficientes de $x^2$ temos
\[
-\left( \frac{1}{ \pi^2}
+ \frac{1}{4\pi^2}
+ \frac{1}{9\pi^2}
+ \cdots
\right)
= - \frac{1}{\,\pi^2\,} \sum_{n=1}^\infty \frac{1}{n^2}    \;.
\]
Igualando os coeficientes temos que
\[
- \frac{1}{6}
= - \frac{1}{\,\pi^2\,} \sum_{n=1}^\infty \frac{1}{n^2}    \;.
\]
Isolando a série chegamos ao resultado desejado
\[
\zeta(2) = \sum_{n=1}^\infty \frac{1}{n^2} = \frac{\pi^2}{6}    \;.
\]

%------------------------------------------------------------------------------%

Apresentamos agora a \emph{demostração alternativa} proposta por Apostol,
que se baseia na avaliação da integral
\[
I = \int_0^1\!\!\int_0^1 \frac{1}{\,1-xy\,}\, dxdy    \;.
\]
O primeiro passo é verificar que essa integral é igual a $\zeta(2)$.
Note que a fração é a soma da Série Geométrica com $\alpha=1$ e $r=xy$,
como os valores de $x$ e $y$ na integral estão entre zero e um
podemos escrever
\[
I = \int_0^1\!\!\int_0^1 \frac{1}{\,1-xy\,} \,dxdy
\;=\; \int_0^1\!\!\int_0^1 \sum_{n=0}^\infty x^n y^n \,dxdy    \;.
\]
Como essa série é absolutamente convergente podemos
trocar a ordem da soma com a integração
\begin{align*}
  I & = \sum_{n=0}^\infty \;\int_0^1\!\!\int_0^1 x^n y^n \,dxdy \\[3mm]
  & = \sum_{n=0}^\infty \;\int_0^1 \frac{y^n}{n+1} \,dy       \\[3mm]
  & = \sum_{n=0}^\infty \frac{1}{(n+1)^2}
  = \zeta(2)    \;.
\end{align*}
A segunda parte consiste em calcular a integral e mostrar que
\[
I = \frac{\pi^2}{6}    \;.
\]
Aplicamos uma mudança de variáveis rotacionando os eixos
\ang{45} no sentido horário
\[
x = \frac{\,u-v\,}{\sqrt{2}} \qquad\qquad
y = \frac{\,u+v\,}{\sqrt{2}}
\]
de modo que
\[
\frac{1}{1-xy} = \frac{2}{\,2-u^2+v^2\,}
\]
e a nova região de integração seja o quadrado, $\Omega$,
com vértices opostos em $(0, 0)$ e $\left(\sqrt{2}, 0\right)$.
O Jacobiano dessa transformação é
\[
J = \frac{\,\partial x\,}{\partial u} \frac{\,\partial y\,}{\partial v}
- \frac{\,\partial x\,}{\partial v} \frac{\,\partial y\,}{\partial u}
= \frac{ 1}{\,\sqrt{2}\,} \frac{1}{\,\sqrt{2}\,}
- \frac{-1}{\,\sqrt{2}\,} \frac{1}{\,\sqrt{2}\,}
= 1    \;.
\]
Assim a integral assume a forma
\[
I = \iint_\Omega \frac{2}{\,2-u^2+v^2\,} \,dudv    \;.
\]
Usando a simetria de reflexão do quadrado em relação
ao eixo $u$ podemos escrever a integral como
\[
I = 4 \int_{0}^{\frac{1}{\sqrt{2}}}
\int_{0}^{u}
\frac{1}{\,2-u^2+v^2\,} \,dvdu
+ 4 \int_{\frac{1}{\sqrt{2}}}^{\sqrt{2}}
\int_{0}^{\sqrt{2}-u}
\frac{1}{\,2-u^2+v^2\,} \,dvdu    \;.
\]
Usando
\[
\int_0^x \frac{dt}{a^2+t^2} = \frac{1}{\,a\,} \arctan\left(\frac{\,x\,}{a}\right)
\]
temos que
\begin{align*}
  \int_{0}^{u} \frac{dv}{\,2-u^2+v^2\,}
  & = \frac{1}{\,\sqrt{2-u^2}\,} \arctan\left(\frac{u}{\,\sqrt{2-u^2}\,}\right) \;, \\[3mm]
  \int_{0}^{\sqrt{2}-u} \frac{dv}{\,2-u^2+v^2\,}
  & = \frac{1}{\,\sqrt{2-u^2}\,} \arctan\left(\frac{\sqrt{2}-u}{\,\sqrt{2-u^2}\,}\right) \;.
\end{align*}
Voltando para o cálculo de $I$ fazemos $I=I_1+I_2$ onde
\begin{align*}
  I_1 & = 4 \int_{0}^{\frac{1}{\sqrt{2}}}
  \arctan\left(\frac{u}{\,\sqrt{2-u^2}\,}\right)
  \frac{du}{\,\sqrt{2-u^2}\,}
  \\[3mm]
  I_2 & = 4 \int_{\frac{1}{\sqrt{2}}}^{\sqrt{2}}
  \arctan\left(\frac{\sqrt{2}-u}{\,\sqrt{2-u^2}\,}\right)
  \frac{du}{\,\sqrt{2-u^2}\,}
\end{align*}
Para calcularmos $I_1$ usamos a transformação
\[
u = \sqrt{2} \sin\left(\theta\right)
\qquad\qquad
du = \sqrt{2} \cos\left(\theta\right) \,d\theta
= \sqrt{2-u^2}\,d\theta
\]
que produz
\begin{align*}
  I_1 & = 4 \int_{0}^{\frac{\pi}{6}}
  \arctan\left(\dfrac{  \sqrt{2} \sin\left(\theta\right)}
  {\,\sqrt{2} \cos\left(\theta\right)\,}\right)
  \,d\theta
  \\[3mm]
  & = 4 \int_{0}^{\frac{\pi}{6}} \theta \,d\theta
  = 2 \left(\frac{\pi}{6}\right)^2
  = \frac{\,2 \pi^2}{6^2}    \;.
\end{align*}
Para $I_2$ usamos \(u = \sqrt{2} \cos\left(2\theta\right)\) que produz
\begin{align*}
  du & = -2\sqrt{2} \sin\left(2\theta\right)                \,d\theta \\[1mm]
  & = -2\sqrt{2} \sqrt{1 - \cos^2\left(2\theta\right)\,} \,d\theta \\[1mm]
  & = -2\sqrt{2} \sqrt{1 - \sfrac{u^2}{2}\,}             \,d\theta \\[1mm]
  & = -2\sqrt{2 - u^2\,}                                 \,d\theta
\end{align*}
e também
\begin{align*}
  \frac{\sqrt{2}-u}{\sqrt{2-u^2\,}}
  & = \frac{\sqrt{2}(1-\cos(2\theta))}{\sqrt{2 - 2\cos^2(2\theta)}} \\[3mm]
  & = \sqrt{\frac{1-\cos(2\theta)}{1 + \cos(2\theta)}\,}            \\[3mm]
  & = \sqrt{\frac{2\sin^2(\theta)}{\cos^2(\theta)}\,}               \\[3mm]
  & = \tan(\theta)  \;.
\end{align*}
Temos então
\[
I_2   =   4 \int_{\frac{\pi}{6}}^{0} \arctan\left(\tan(\theta)\right) (-2)\,d\theta
\;=\; 8 \int_{0}^{\frac{\pi}{6}} \theta \,d\theta
\;=\; 4 \left(\frac{\pi}{6}\right)^2
\;=\; \frac{\,4 \pi^2}{6^2}  \;.
\]
Calculando $I$ temos
\[
I = I_1 + I_2
= \frac{\,2 \pi^2}{6^2}
+ \frac{\,4 \pi^2}{6^2}
= \frac{\pi^2}{6}
\]
o que conclui nossa demostração.

